\section{Results}
\subsection{RQ1: Does Unlabeled Adaptation Improve P0?}
T3A improves receiver-equal macro-F1 from \NPZeroMacroFOne{} to \NTThreeAMacroFOne{}. The absolute paired improvement is \TGain{}, positive on 31 of the 32 receivers. The stored receiver-bootstrap interval is [\TLower{}, \TUpper{}]. The two-sided Monte Carlo sign-flip value is approximately $10^{-5}$; its Holm-adjusted value is unchanged in the later benchmark's one-comparison family. These statistics support consistency across the sampled receivers, subject to the evidence-timing and shared-training limitations in Section III.

Table~\ref{tab:bench} reports all primary benchmark entries, including weaker outcomes. T3A is the strongest evaluated deployable same-information method. This is a descriptive ranking within the frozen candidate set, not state-of-the-art superiority over all RF or adaptation methods.
\begin{table}[t]
\centering
\caption{Primary receiver-equal macro-F1. SD is across receivers after seed averaging, not a confidence interval.}
\label{tab:bench}
\small
\begin{tabular}{lrrr}
\toprule
Method & Mean & Receiver SD & ECE\\
\midrule
P0 & 0.805679 & 0.0843 & 0.0889\\
P0\_WIDE & 0.801535 & 0.0900 & 0.0997\\
CORAL & 0.805071 & 0.0826 & 0.0853\\
DANN & 0.800926 & 0.0848 & 0.0829\\
GROUPDRO & 0.761754 & 0.0789 & 0.0844\\
RX\_NORM & 0.800769 & 0.0869 & 0.0919\\
T3A & 0.833692 & 0.0789 & 0.1030\\
P1 & 0.809384 & 0.0860 & 0.1039\\
P2 & 0.806726 & 0.0820 & 0.1083\\
SUP-FT oracle & 0.838081 & 0.0814 & 0.0648\\
\bottomrule
\end{tabular}
\end{table}

\begin{figure*}[t]\centering\includegraphics[width=.95\textwidth]{fig02_benchmark}
\caption{Receiver-equal macro-F1 for all benchmark methods. Error bars are receiver standard deviations after averaging seeds, not inferential confidence intervals.}\label{fig:benchmark}\end{figure*}
Figure~\ref{fig:benchmark} emphasizes both the ranking and substantial spread across receivers. Absolute performance remains receiver dependent even when a method improves most paired comparisons.

\subsection{RQ2: Which Evaluated Strategy Works Best?}
P2 reaches \NPTwoMacroFOne{}, a mean difference of only \PGain{} from P0. Its P2--T3A difference is \PvsT{}. P0-WIDE, the source-only capacity control, reaches \NPZeroWIDEMacroFOne{}. The source-only CORAL, DANN, and GroupDRO variants do not exceed P0's mean. This does not establish that source-domain robustness objectives are ineffective in general; they were evaluated under a particular backbone, budget, and domain definition.

The supervised fine-tuning oracle reaches \NSUPFTOneTwoEightMacroFOne{}. It uses the same number of receiver packets but receives support labels, so it must remain outside an unlabeled-method superiority claim. Its modest mean separation from T3A is informative about this particular head-adaptation procedure, not a bound on all possible supervised adaptations.
\begin{table}[t]
\centering
\caption{Paired receiver differences after five-seed averaging. Only T3A--P0 belongs to the later benchmark's one-test family; other entries are descriptive.}
\label{tab:paired}
\small
\begin{tabular}{lrr}
\toprule
Comparison & Mean delta & Positive / 32\\
\midrule
T3A - P0 & +0.028014 & 31\\
P2 - P0 & +0.001047 & 15\\
P2 - T3A & -0.026966 & 5\\
P2 - P0\_WIDE & +0.005191 & 20\\
P2 - P2\_SHUFFLED & +0.018364 & 32\\
P2 - P2\_MISMATCHED\_RX & +0.019637 & 30\\
\bottomrule
\end{tabular}
\end{table}


\begin{figure*}[t]\centering\includegraphics[width=.95\textwidth]{fig03_receivers}
\caption{T3A--P0 differences for every physical receiver after five-seed averaging. The symmetric difference axis includes zero; receivers are not sorted to conceal negative outcomes.}\label{fig:receivers}\end{figure*}
\subsection{RQ3: How Much Support Is Needed?}
The common-query diagnostic in Table~\ref{tab:budget} reveals strong support dependence. T3A is substantially worse at the smallest bank and improves as the bank grows. P2 varies much less across the prescribed budgets and does not obtain a corresponding performance gain. No best bank size is selected: the original primary budget remains fixed.
\begin{figure*}[t]\centering\includegraphics[width=.95\textwidth]{fig04_budget}
\caption{Full post-hoc budget curves. Lines connect measured support sizes, not acquisition time. This diagnostic uses a query pool with the maximum bank reserved.}\label{fig:budget}\end{figure*}
Figure~\ref{fig:budget} should not be read as evidence that any unlabeled information is automatically helpful. Prototype updates can be harmful when bounded support is insufficient. The diagnostic is not numerically identical to the primary table, because query membership differs.
\begin{table}[t]
\centering
\caption{Post-hoc common-query support-budget diagnostic. The pool reserves 256 packets; these means are not the primary query-set means.}
\label{tab:budget}
\small
\begin{tabular}{lrrr}
\toprule
Support & T3A & P2 & RX-NORM\\
\midrule
16 & 0.7196 & 0.8035 & 0.7701\\
32 & 0.7956 & 0.8063 & 0.7882\\
64 & 0.8228 & 0.8066 & 0.7957\\
128 & 0.8336 & 0.8067 & 0.8008\\
256 & 0.8383 & 0.8068 & 0.8033\\
\bottomrule
\end{tabular}
\end{table}


\subsection{RQ4: How Heterogeneous Are Receiver Effects?}
The paired view in Fig.~\ref{fig:receivers} retains all receivers in identifier order. T3A's positive mean is not generated by a small favorable subset, although one receiver has a negative seed-averaged difference. P2's positive and negative receiver counts are much more balanced. Hardware-family summaries likewise report positive and negative context effects rather than selecting a favorable family.

These patterns motivate receiver-level reporting. Repeated packets from one receiver cannot supply the same hardware-transfer evidence as additional independent receivers. Nonetheless, LOSO training sets overlap, so the receiver bootstrap remains conditional on this dataset and procedure.

\subsection{RQ5: Does Classification Improvement Improve Calibration?}
It does not uniformly do so. T3A's ECE is \NTThreeAEce{}, compared with P0's \NPZeroEce{}, despite higher accuracy. P2 has ECE \NPTwoEce{}. T3A's NLL is lower than P0's, showing that confidence calibration cannot be reduced to a single scalar judgment. No temperature or confidence threshold is fitted to improve these results.
\begin{figure*}[t]\centering\includegraphics[width=.95\textwidth]{fig05_calibration}
\caption{Classification and probability diagnostics. Left: every receiver for P0, T3A, and P2. Right: ECE and NLL are distinct measures, not interchangeable units.}\label{fig:calibration}\end{figure*}
Figure~\ref{fig:calibration} places accuracy--ECE relationships next to the distinct score measures. T3A changes the classifier's scoring geometry; the tradeoff matters when downstream use needs trustworthy probabilities rather than only a class decision.

\subsection{RQ6: What Are the Computational Costs?}
T3A reuses P0's trained encoder and adjusts prototypes without gradients. P2 requires separate source training and repeated support encoding or a compatible embedding cache. Table~\ref{tab:compute} distinguishes these boundaries. Logged adaptation/evaluation durations include operations beyond adaptation itself and do not establish isolated per-query deployment latency.
\begin{figure*}[t]\centering\includegraphics[width=.95\textwidth]{fig06_compute}
\caption{Frozen cost accounting. Source training and reused-checkpoint adaptation/evaluation occupy separate panels because timing scopes differ.}\label{fig:compute}\end{figure*}
T3A's practical appeal is a bounded conclusion: it improves this benchmark without an additional source-trained context module. Figure~\ref{fig:compute} is not a hardware-independent throughput benchmark. Missing timing components are disclosed, not imputed.
\begin{table}[t]
\centering
\caption{Logged cost boundaries differ. Adapt/eval records reuse source checkpoints; timings are not isolated per-packet latency.}
\label{tab:compute}
\small
\begin{tabular}{lrrr}
\toprule
Method & Parameters & Seconds & Scope\\
\midrule
P0 & 64774 & 30.992 & Source train\\
P2 & 75143 & 98.296 & Source train\\
T3A & 64774 & 0.671 & Adapt/eval\\
RX-NORM & 64774 & 0.492 & Adapt/eval\\
SUP-FT & 64774 & 0.689 & Adapt/eval\\
\bottomrule
\end{tabular}
\end{table}


\subsection{RQ7: What Did Learned Context Conditioning Fail to Capture?}
P2 performs better with natural receiver support than with shuffled or mismatched donors, but this sensitivity is insufficient to establish useful adaptation. The natural-support mean remains approximately equal to P0 and below T3A. A model can depend on context without using it better than a competent independent or prototype-adjusted comparator.

Post-hoc inference-only query-coupling diagnostics do not materially recover the initial context study's apparent advantage. Shuffled-context source training retains sensitivity to matched support but yields an absolute mean below P0. Oracle composition controls alter multiple methods, including T3A and receiver normalization; they do not isolate a uniquely P2-specific explanation. The equalized matched-ID diagnostic was ineligible under frozen identifiers. These findings are post-hoc and cannot replace the primary V2 decision.
